\documentclass[11pt,a4paper]{article}
\usepackage{microtype}
\usepackage[margin=2.6cm]{geometry}
\usepackage{amsmath,amsthm,thmtools,thm-restate}
\usepackage{fontspec}
\usepackage{unicode-math}
\directlua{luaotfload.add_fallback("cfb", {"NotoSansMath:mode=harf;", "DejaVuSans:mode=harf;"})}
\setmainfont{Libertinus Serif}[RawFeature={fallback=cfb}]
\setmathfont{Latin Modern Math}
\setmonofont{Latin Modern Mono}[Scale=0.85,RawFeature={fallback=cfb}]
\AtBeginDocument{\let\setminus\smallsetminus}
\usepackage{enumitem}
\usepackage{longtable,booktabs,array,calc}
\usepackage{tcolorbox}
\usepackage{xurl}
\usepackage{fancyhdr}
\usepackage[colorlinks=true,linkcolor=blue!50!black,citecolor=blue!50!black,urlcolor=blue!50!black]{hyperref}
\usepackage[capitalize,nameinlink]{cleveref}
\setlength{\emergencystretch}{3em}

\theoremstyle{plain}
\newtheorem{theorem}{Theorem}[section]
\newtheorem{lemma}[theorem]{Lemma}
\newtheorem{corollary}[theorem]{Corollary}
\newtheorem{proposition}[theorem]{Proposition}
\newtheorem{observation}[theorem]{Observation}
\theoremstyle{definition}
\newtheorem{definition}[theorem]{Definition}
\newtheorem{question}[theorem]{Question}
\newtheorem{conjecture}[theorem]{Conjecture}
\theoremstyle{remark}
\newtheorem{remark}[theorem]{Remark}
\newtheorem*{proofidea}{Idea of proof}
\newcommand{\fullproofref}[1]{\,\hyperref[#1]{\textit{[full proof]}}}
\providecommand{\tightlist}{\setlength{\itemsep}{0pt}\setlength{\parskip}{0pt}}

\newcommand{\E}{\mathbb E}
\newcommand{\R}{\mathbb R}
\newcommand{\one}{\mathbf 1}
\newcommand{\cT}{\mathcal T}
\newcommand{\cO}{\mathcal O}
\newcommand{\per}{\operatorname{per}}
\newcommand{\tr}{\operatorname{tr}}
\newcommand{\repair}[1]{\begin{tcolorbox}[colback=yellow!7,colframe=gray!60,boxrule=0.3pt,arc=1pt,left=5pt,right=5pt,top=3pt,bottom=3pt]\small\textbf{Repair.} #1\end{tcolorbox}}

\pagestyle{fancy}\fancyhf{}
\fancyhead[L]{\small\itshape Candidate result chromatic\_number\_of\_random\_lifts\_of\_complete\_graphs --- machine-generated proof, awaiting human review}
\fancyhead[R]{\small\thepage}
\renewcommand{\headrulewidth}{0.2pt}

\title{Random lifts of $K_5$ are asymptotically almost surely $3$-chromatic}
\author{\href{https://github.com/graph-theory-AI}{github.com/graph-theory-AI}\\[2pt] \normalsize Machine-generated note}
\date{18 September 2026}

\begin{document}
\maketitle

\begin{abstract}
We prove that a uniform random $n$-lift of $K_5$ has chromatic number $3$ with
probability tending to one.  A Bregman permanent bound keeps the matching constraints
in the second moment of fibre-balanced $3$-colourings, making the uniform overlap the
unique global maximum; the resulting variance is then exhausted by short cycles, so
small-subgraph conditioning yields a colouring asymptotically almost surely.

\medskip\noindent\small This note was produced by AI within
the \href{https://github.com/graph-theory-AI/Graph-Theory-LLM-Proofs}{Graph Theory AI}
project: the argument was found by a language model and audited by an adversarial LLM
referee, and it has not been checked by a human mathematician. \Cref{sec:provenance}
records the full provenance.
\end{abstract}

\section{Introduction}

An $n$-lift of a graph $G$ replaces every base vertex by a fibre of $n$ vertices and
every base edge by a perfect matching between the corresponding fibres.  In a uniform
random lift these matchings are chosen independently and uniformly.  The
OpenProblemGarden entry \cite{OPG}, following Amit, Linial and Matoušek \cite{ALM02},
asks verbatim:
\begin{question}\label{q:opg}
Is the chromatic number of a random lift of $K_5$ concentrated on a single value?
\end{question}
The possible asymptotic values are $3$ and $4$.  Farzad and Theis proved the analogous
$3$-colourability result for $K_5\setminus e$ \cite{FT12}.

Nir and Pérez-Giménez \cite{NPG21} locate the chromatic number of random lifts of
$K_{d+1}$ in a two-value window.  For three colours their upper theorem requires
$d<\ell_3=16\log2/3\approx3.69678$, while their lower theorem starts at
$u_3=2\log3/(\log3-\log2)\approx5.41902$.  Thus $d=4$, corresponding to $K_5$, lies
in the gap.  Their second-moment optimization relaxes away matching marginals.  The
argument here instead uses Bregman's permanent inequality \cite{Bregman73}, in its
entropy form due to Radhakrishnan \cite{Radhakrishnan97}, and remains sharp at the
uniform overlap when $d=4$.

\begin{restatable}[Chromatic number of a random lift of $K_5$]{theorem}{thmmain}\label{thm:main}
Let $H_n$ be a uniform random $n$-lift of $K_5$.  Then
\[
 \Pr\bigl(\chi(H_n)=3\bigr)\longrightarrow1.
\]
\end{restatable}
\begin{proofidea}
Count colourings that use nearly $n/3$ vertices of each colour in every fibre.  An
entropy inequality and a permanent bound show that only nearly independent pairs of
colourings contribute to the second moment.  A local lattice calculation evaluates
that contribution exactly.  Its excess over the squared first moment equals precisely
the contribution of short cycles, by the Ihara--Bass determinant identity.  The small-
subgraph conditioning theorem then shows that a balanced $3$-colouring exists with
probability tending to one.  A first-moment estimate excludes bipartite lifts.
\fullproofref{pf:main}
\end{proofidea}

\section{Balanced colourings and the two analytic inputs}

Choose integers $t_1,t_2,t_3$ summing to $n$, each differing from $n/3$ by at most
$1$, and set $p_a=t_a/n$.  Let $Z_n$ count proper $3$-colourings in which every fibre
has exactly $t_a$ vertices of colour $a$.  Write
\[
 M_n=\frac{n!}{t_1!t_2!t_3!},
\]
and let $a_n$ be the probability that a uniform perfect matching between two fibres
with these colour counts has no monochromatic edge.  Then
\[
 \E Z_n=M_n^5a_n^{10}.                                         \tag{1}
\]

For two balanced colourings, the overlap in fibre $i$ is the matrix
\[
 \rho_i(a,b)=\frac1n\bigl|\{x:\text{the two colours of $x$ are $a,b$}\}\bigr|.
\]
Its row and column sums are $p$.  Independently in the five fibres,
\[
 \Pr(\rho_i=\rho)=
 \frac{\prod_{a=1}^3(t_a!)^2}{n!\prod_{a,b=1}^3(n\rho(a,b))!}.  \tag{2}
\]
If $q_n(\rho,\sigma)$ is the probability that a matching between fibres of overlap
types $\rho,\sigma$ respects both colourings, then
\[
 \frac{\E Z_n^2}{(\E Z_n)^2}
 =\E_{\rho_1,\ldots,\rho_5}
 \prod_{ij\in E(K_5)}\frac{q_n(\rho_i,\rho_j)}{a_n^2}.           \tag{3}
\]

Let $U$ be the $3\times3$ matrix with all entries $1/9$, and define
\[
 H(\rho)=-\sum_{a,b}\rho(a,b)\log\rho(a,b),\qquad
 L(\rho)=\sum_{a,b}\rho(a,b)\log\bigl(\tfrac13+\rho(a,b)\bigr),
\]
and $F=H+2L$.

\begin{restatable}[Entropy inequality]{lemma}{lementropy}\label{lem:entropy}
There is $c>0$ such that every nonnegative $3\times3$ matrix with row sums $1/3$
satisfies
\[
 F(\rho)\le F(U)-c\|\rho-U\|_F^2,
 \qquad F(U)=2\log(4/3).
\]
\end{restatable}
\begin{proofidea}
After rescaling a row by $3$, the problem is to maximize
$\sum_{j=1}^3g(x_j)$ on the simplex for
$g(x)=2x\log(1+x)-x\log x$.  The function has a single inflection point.  Concavity
handles three small coordinates, while reducing the one-large-coordinate case to a
strictly concave one-variable function.  The Hessian at the uniform point is negative
definite, giving the quadratic gap.\fullproofref{pf:entropy}
\end{proofidea}

\begin{restatable}[Finite-type matching asymptotic]{lemma}{lemmatching}\label{lem:matching}
Let $A$ be a symmetric zero-one matrix of order $s$ with constant row sum $r$, put
$P=A/r$ and $u=(1/s,\ldots,1/s)$.  Assume that $I-P^2$ is positive definite on
$\one^\perp$ and that the vectors $e_j-e_k$ for columns occurring together in a row
support generate the zero-sum integer lattice.  If $Q_A(n;\alpha,\beta)$ is the
probability that a uniform matching with integer type counts $n\alpha,n\beta$ uses only
allowed type pairs, then, uniformly for
$\|\alpha-u\|+\|\beta-u\|=O(n^{-2/5})$,
\begin{align*}
Q_A(n;\alpha,\beta)
={}&\left(\frac rs\right)^n
 \det_{\one^\perp}(I-P^2)^{-1/2}\\
&\times\exp\left\{\frac{ns}{2}
 \left[\|\beta-u\|^2-
 \langle w,(I-P^2)^{-1}w\rangle\right]\right\}(1+o(1)),\\
&\hspace{35mm}w=\beta-u-P(\alpha-u).
\end{align*}
\end{restatable}
\begin{proofidea}
Choose permitted right types independently and express the matching probability as a
ratio of two lattice point probabilities.  Exponential tilting gives the quadratic rate
function; uniform Fourier inversion supplies the Gaussian prefactor.  The lattice
hypothesis rules out nontrivial maximum-modulus characters.\fullproofref{pf:matching}
\end{proofidea}

\section{Moment and cycle calculations}

\begin{restatable}[Second moment]{proposition}{propsecond}\label{prop:second}
For the variable $Z_n$ above,
\[
 \frac{\E Z_n^2}{(\E Z_n)^2}\longrightarrow
 R:=25\left(\frac{16}{15}\right)^{20}\left(\frac{15}{23}\right)^8.
\]
Moreover $\E Z_n=\Theta(n^{-5})(4/3)^{5n}$.
\end{restatable}
\begin{proofidea}
Bregman's inequality and \cref{lem:entropy} suppress every overlap farther than
$n^{-2/5}$ from independence.  On the central window, \cref{lem:matching} and Stirling
turn the sum into a $20$-dimensional Gaussian integral.  Its five-fibre matrix is
$(19I-4A(K_5))/15$, with eigenvalues $1/5,23/15,\ldots,23/15$.
\fullproofref{pf:second}
\end{proofidea}

Let $B$ be the nonbacktracking matrix of $K_5$, indexed by its $20$ directed edges,
and let $X_{\ell,n}$ count unoriented simple $\ell$-cycles in $H_n$.  Put
\[
 \mu_\ell=\frac{\tr(B^\ell)}{2\ell},\qquad
 \delta_\ell=2(-1/2)^\ell.                                    \tag{4}
\]

\begin{restatable}[Short-cycle moments and variance identity]{proposition}{propcycles}\label{prop:cycles}
For every fixed $L$, $(X_{3,n},\ldots,X_{L,n})$ converges jointly to independent
Poisson variables of means $(\mu_3,\ldots,\mu_L)$.  For every finitely supported
sequence of nonnegative integers $b_\ell$,
\[
 \frac{\E\left[Z_n\prod_{\ell\ge3}(X_{\ell,n})_{b_\ell}\right]}{\E Z_n}
 \longrightarrow
 \prod_{\ell\ge3}\bigl(\mu_\ell(1+\delta_\ell)\bigr)^{b_\ell}.
\]
Finally,
\[
 \sum_{\ell\ge3}\mu_\ell\delta_\ell^2<\infty,
 \qquad
 R=\exp\left(\sum_{\ell\ge3}\mu_\ell\delta_\ell^2\right).
\]
\end{restatable}
\begin{proofidea}
Factorial moments count vertex-disjoint lifted cycles; intersecting configurations lose
a factor $n^{-1}$.  Under the colouring-size-biased measure, a cycle gains the factor
$1+\delta_\ell$.  The Ihara--Bass determinant identity at $u=1/4$ turns the sum of
$\mu_\ell\delta_\ell^2$ into the same determinant as the Gaussian second moment.
\fullproofref{pf:cycles}
\end{proofidea}

\section{Remarks}

\begin{remark}[What is new relative to the 2021 bound]
The relaxation used in \cite{NPG21} has a nonuniform maximizer when $d=4$.  The
permanent bound here retains the perfect-matching constraint and has the uniform
overlap as its unique maximizer.  The referee numerically optimized the exact
maximum-entropy coupling as an independent check and found the same maximizer.
\end{remark}

\begin{remark}[Scope]
The argument also gives asymptotic $3$-colourability for lifts of any fixed simple
$4$-regular base; it does not assert chromatic number exactly $3$ for bipartite bases.
The general OpenProblemGarden question for $K_m$ remains untouched.  The entropy lemma
is specific to degree $4$: its analogue already fails at degree $5$, so this method
does not resolve random lifts of $K_6$.
\end{remark}

\begin{remark}[Computational audit]
The referee verified every displayed constant.  In particular,
$R=2.9745540512$ and
\[
 \sum_{\ell=3}^{60}\mu_\ell\delta_\ell^2=1.090094126,
 \qquad \log R=1.090094129,
\]
the residual being truncation.  In simulations, the $3$-colourable fraction was
$211/400$ at lift size $3$, $293/300$ at size $12$, and $200/200$ at size $20$;
no bipartite lift occurred in more than $1{,}400$ trials.  These computations support
the proof but do not replace it.
\end{remark}

\section{Provenance}\label{sec:provenance}

The mathematics in this note was produced without human intervention, in three stages.
(1)~The proof was found by \texttt{gpt-6-astra}, reasoning effort \texttt{max}, in a
single pass started on 2026-09-16, for OpenProblemGarden catalog id
\texttt{chromatic\_number\_of\_random\_lifts\_of\_complete\_graphs}, campaign leg
\texttt{attacks\_opg}.
(2)~An independent adversarial referee agent (\texttt{claude-opus-5}) re-derived the
entropy, permanent, Gaussian, cycle and determinant calculations, checked the literature,
and ran the computations summarized above on 2026-09-17; its verdict was
\textsc{minor gaps}.  Its report is reproduced verbatim in \cref{app:referee}.
(3)~No rewrite of this target was started by \texttt{claude-fable-5-1} before its
session ended; only the original writeup and referee report were present.
(4)~On 2026-09-18 \texttt{Codex} wrote the present note and moved the full proofs to
\cref{app:proofs}.  It repaired and marked four referee gaps: a finite-dimensional
entropy-continuity proof for rounded fibre sizes; a uniform Fourier argument for the
local lattice estimate, with \cite{GJR10}; a full factorial-moment enumeration explaining
periodic projected walks; and an explicit invocation of the small-subgraph conditioning
theorem from \cite{Wormald99}, with every hypothesis matched to a preceding formula.
Citations for Bregman's theorem and Radhakrishnan's entropy proof were also added.  No
other mathematical content was changed.  \textbf{No human mathematician has checked
the argument.}

\appendix
\section{Full proofs}\label{app:proofs}

\phantomsection\label{pf:entropy}
\lementropy*
\begin{proof}
Put $g(x)=2x\log(1+x)-x\log x$ for $0\le x\le1$, with $0\log0=0$.
We first prove
\[
 g(x_1)+g(x_2)+g(x_3)\le3g(1/3)\qquad(x_1+x_2+x_3=1),           \tag{5}
\]
with equality only at the uniform point.  Direct differentiation gives
\[
 g''(x)=\frac{x^2+2x-1}{x(1+x)^2}.
\]
Thus $g$ is strictly concave below $r=\sqrt2-1$ and strictly convex above it.  A
maximizer is interior because $g'(0+)=+\infty$.  At an interior maximizer, two
coordinates cannot both exceed $r$, since varying those two with fixed sum has positive
second derivative.

If every coordinate is at most $r$, Jensen's inequality proves (5).  Otherwise, write
the exceptional coordinate as $x>r$.  Concavity on the other two gives
\[
 g(x_1)+g(x_2)+g(x_3)\le \phi(x):=g(x)+2g((1-x)/2).
\]
For $1/3\le x<1$,
\[
 \phi''(x)=
 \frac{x^4-10x^3-32x^2+34x-9}{x(1+x)^2(1-x)(3-x)^2}<0.
\]
Indeed, the numerator is at most $-(103/3)x^2+34x-9$, whose discriminant is $-80$.
Since $\phi'(1/3)=0$, strict concavity gives $\phi(x)<\phi(1/3)$ for $x>1/3$.

Apply (5) to each row of $x_{ab}=3\rho(a,b)$.  Since
\[
 F(\rho)=-\log3+\frac13\sum_{a,b}g(3\rho(a,b)),
\]
$U$ is the unique maximizer.  The Hessian of $F$ at $U$ is $-(9/8)I$ in the nine
entry coordinates.  Taylor's theorem gives a strict quadratic decrease near $U$;
on the compact complement of a small neighbourhood, uniqueness gives a fixed gap,
which may be weakened to the required quadratic bound.
\end{proof}

\phantomsection\label{pf:matching}
\lemmatching*
\begin{proof}
For every left vertex, independently choose one of its $r$ permitted right types
uniformly, and let $S$ be the resulting vector of right-type counts.  Counting matchings
with a fixed type table in two ways gives the exact identity
\[
 Q_A(n;\alpha,\beta)=
 \left(\frac rs\right)^n
 \frac{\Pr(S=n\beta)}{\Pr(\operatorname{Mult}(n,u)=n\beta)}.    \tag{6}
\]
The mean of $S/n$ is $P\alpha$.  At $\alpha=u$, its covariance divided by $n$, on
$\one^\perp$, is
\[
 \Sigma_0=\frac1s(I-P^2).
\]
The normalized logarithmic moment generating function is
\[
 K_\alpha(t)=\sum_i\alpha_i\log\left(\sum_jP_{ij}e^{t_j}\right).
\]
Uniformly near $(u,0)$ on $\one^\perp$,
\[
 K_\alpha(t)=\langle P\alpha,t\rangle+	frac12\langle
 t,\Sigma_0t\rangle+cO(\|\alpha-u\|\|t\|^2+\|t\|^3).
\]
Writing $w=\beta-u-P(\alpha-u)$, the tilt with mean $\beta$ is
$t=\Sigma_0^{-1}w+cO((\|\alpha-u\|+\|\beta-u\|)^2)$, and its rate function is
\[
 I_\alpha(\beta)=\tfrac12\langle w,\Sigma_0^{-1}w\rangle
 +\cO((\|\alpha-u\|+\|\beta-u\|)^3).                           \tag{7}
\]

\repair{The original writeup asserted uniform Fourier inversion.  Here are the
uniformity checks.  After tilting, the type proportions remain in a fixed compact
neighbourhood of $u$ and the covariance on $\one^\perp$ remains uniformly positive
definite.  The characteristic function is a product of finitely many categorical
characteristic functions.  The lattice-generation hypothesis says that equality in
the triangle inequality on the dual torus occurs only at the trivial character;
compactness therefore gives modulus at most $e^{-cn}$ outside a fixed neighbourhood
of zero, uniformly in $\alpha,\beta$.  Inside that neighbourhood, the logarithm has
the uniform expansion $-n\langle t,\Sigma t\rangle/2+cO(n\|t\|^3)$.  Splitting at
$\|t\|=n^{-2/5}$ makes the cubic error $o(1)$ on the Gaussian scale, the intermediate
annulus is bounded by $e^{-cn^{1/5}}$, and rescaling $t$ by $n^{-1/2}$ gives the
Gaussian integral uniformly.  This is the fixed-dimensional lattice Laplace argument
of Greenhill--Janson--Ruciński \cite{GJR10}; all its nondegeneracy and lattice
hypotheses are exactly the two assumptions of this lemma.}

Consequently, the local probability in the numerator of (6) has the Gaussian
prefactor on the full zero-sum lattice.  The denominator has the same lattice
covolume, so it cancels.  The determinant ratio is
$\det_{\one^\perp}(I-P^2)^{-1/2}$.  Finally,
\[
 D(\beta\|u)=\frac s2\|\beta-u\|^2+cO(\|\beta-u\|^3).
\]
The deviations are $O(n^{-2/5})$, so multiplying the cubic errors by $n$ gives
$O(n^{-1/5})=o(1)$.  Substitution in (6), using
$\Sigma_0^{-1}=s(I-P^2)^{-1}$, proves the displayed formula uniformly.
\end{proof}

\phantomsection\label{pf:second}
\propsecond*
\begin{proof}
We first obtain a global bound.  Bregman's theorem states that a zero-one matrix with
positive row sums $r_v$ satisfies
\[
 \per(A)\le\prod_v(r_v!)^{1/r_v}.                               \tag{8}
\]
For the matching counted by $q_n(\rho,\sigma)$, a left vertex of type $(a,b)$ has
\[
 n(1-p_a-p_b+\sigma(a,b))
\]
permitted partners, uniformly at least $n/4$ for large $n$.  Applying (8), Stirling's
formula, and then interchanging $\rho,\sigma$ gives
\begin{align*}
 \log q_n(\rho,\sigma)
 &\le n\sum_{a,b}\rho(a,b)\log(\tfrac13+\sigma(a,b))+cO(\log n),\\
 \log q_n(\rho,\sigma)
 &\le \frac n2\bigl(L(\rho)+L(\sigma)\bigr)+\cO(\log n).       \tag{9}
\end{align*}
For the second line, use monotonicity of the logarithm to obtain
\[
 \sum_{a,b}(\rho(a,b)-\sigma(a,b))
 [\log(\tfrac13+\rho(a,b))-\log(\tfrac13+\sigma(a,b))]\ge0.
\]

\repair{For the rounded margins, rescale row $a$ of $\rho$ by $1/(3p_a)$, obtaining
$\widehat\rho$ with row sums $1/3$ and
$\|\rho-\widehat\rho\|_1=O(1/n)$.  In fixed dimension, the elementary entropy
continuity bound
$|H(x)-H(y)|\le h_2(\delta)+\delta\log(8)=O(\delta\log(1/\delta))$ for
$\delta=\|x-y\|_1/2$ gives an $O((\log n)/n)$ change in $H$.  The derivative of
$x\mapsto x\log(1/3+x)$ is bounded on $[0,1]$, so $L$ changes by $O(1/n)$.  Since
$\|\widehat\rho-U\|_F^2\ge\tfrac12\|\rho-U\|_F^2-O(n^{-2})$,
\cref{lem:entropy} yields, uniformly,
\[
 F(\rho)\le F(U)-c'\|\rho-U\|_F^2+O((\log n)/n).                \tag{10}
\]}

Stirling's formula in (2) gives
\[
 \log\Pr(\rho_i=\rho)=n(H(\rho)-2H(p))+O(\log n).
\]
Also $H(p)=\log3+O(n^{-2})$.  Combining this with (9), the first-moment asymptotic
derived below, and the fact that every fibre of $K_5$ has degree $4$, the contribution
of a fixed five-overlap profile to (3) is at most
\[
 \exp\left\{n\sum_{i=1}^5(F(\rho_i)-F(U))+O(\log n)\right\}.    \tag{11}
\]
Put $C_n=pp^{T}$.  Since $C_n-U=O(1/n)$, (10)--(11) show that profiles with some
$\|\rho_i-C_n\|_F>n^{-2/5}$ contribute at most a polynomial times
$e^{-cn^{1/5}}$.  There are only polynomially many profiles, so their total is $o(1)$.

It remains to evaluate the central window.  Apply \cref{lem:matching} first to
$A_0=J_3-I_3$.  The eigenvalues of $P_0=A_0/2$ are $1,-1/2,-1/2$, and the lattice
hypothesis holds, whence
\[
 a_n=\frac43(2/3)^n(1+o(1)).                                  \tag{12}
\]
Together with Stirling's formula for $M_n$, (1) gives
$\E Z_n=\Theta(n^{-5})(4/3)^{5n}$.

Let
\[
 \cT=\{X\in\R^{3\times3}:X\one=X^T\one=0\},\qquad
 z_i=3\sqrt n(\rho_i-C_n).
\]
This is four-dimensional.  Uniformly in the central window, (2) and Stirling give
\[
 \Pr(\rho_i=\rho)=\frac{729}{(2\pi n)^2}e^{-\|z_i\|_F^2/2}(1+o(1)). \tag{13}
\]
The row-and-column-zero integer lattice has covolume $9$: the basis
$(e_a-e_3)(e_b-e_3)^T$, $a,b\in\{1,2\}$, has Gram determinant $81$.  Hence the
$z_i$-lattice has covolume $729/n^2$, matching (13).

For $A_1=A_0\otimes A_0$, $P_1=A_1/4=P_0\otimes P_0$ has eigenvalues
$1$, $-1/2$ with multiplicity $4$, and $1/4$ with multiplicity $4$; on $\cT$ it acts
as multiplication by $1/4$.  The lattice hypothesis again holds.  Applying
\cref{lem:matching} and dividing by (12) squared yields
\[
 \frac{q_n(\rho_i,\rho_j)}{a_n^2}
 =\left(\frac{16}{15}\right)^2
 \exp\left\{\frac4{15}\langle z_i,z_j\rangle
 -\frac1{30}(\|z_i\|^2+\|z_j\|^2)\right\}(1+o(1)).             \tag{14}
\]

Let $A$ now be the adjacency matrix of $K_5$ and set
\[
 M=\frac{19I_5-4A}{15}.
\]
Equations (13)--(14) give this quadratic form independently in each of the four
coordinates of $\cT$.  Its eigenvalues are $1/5$ once and $23/15$ four times, so it is
positive definite.  The central lattice sum therefore converges to its Gaussian
integral, and
\[
 \frac{\E Z_n^2}{(\E Z_n)^2}\longrightarrow
 \left(\frac{16}{15}\right)^{20}\det(M)^{-2}
 =25\left(\frac{16}{15}\right)^{20}\left(\frac{15}{23}\right)^8.
\]
\end{proof}

\phantomsection\label{pf:cycles}
\propcycles*
\begin{proof}
A simple cycle in the lift projects to a cyclically nonbacktracking closed walk in the
base.  For a fixed rooted oriented projected walk of length $\ell$, choosing distinct
lift vertices gives $n^\ell(1+O(1/n))$ possibilities, and the required matching edges
occur with probability $n^{-\ell}(1+O(1/n))$.  The factorial-moment expansion for a
fixed finite collection of lengths is dominated by vertex-disjoint cycles.  Such
choices factor.  If two distinct cycles overlap, their union has strictly more edges
than vertices and contributes $O(n^{v-e})=O(1/n)$.  Dividing by the $2\ell$ roots and
orientations of each lifted cycle gives the independent Poisson means
$\mu_\ell=\tr(B^\ell)/(2\ell)$.

\repair{This rooted count also handles a projected walk that is a proper power.  If a
primitive base walk is traversed $d$ times, closure in the lift asks for a $d$-cycle of
its monodromy permutation, whose expected number is $1/d$.  Equivalently, its periodic
word has only $2\ell/d$ distinct base root-orientations, and the factor $1/d$ restores
the same normalization by $2\ell$.  Thus periodic projections, such as doubled
triangles at $\ell=6$, are already included in $\tr(B^\ell)/(2\ell)$; configurations
with repeated lift vertices belong to the $O(1/n)$ degenerate terms.}

Now bias the lift law by $Z_n/\E Z_n$.  Equivalently, choose a uniform balanced
colouring in each fibre and then independent uniform colour-respecting matchings.  For
$r$ prescribed disjoint compatible edges in one matching, the probability that they
all occur is
\[
 (3/(2n))^r(1+o(1)),                                           \tag{15}
\]
uniformly for fixed $r$.  This follows by deleting the endpoints and applying
\cref{lem:matching} to the remaining counts, which differ from uniform by $O(1/n)$.
The planted colours on a fixed vertex set are asymptotically independent and uniform.
An $\ell$-cycle has $2^\ell+2(-1)^\ell$ proper $3$-colourings, so its relative weight
under the biased law is
\[
 3^{-\ell}(2^\ell+2(-1)^\ell)(3/2)^\ell=1+\delta_\ell.          \tag{16}
\]

For the mixed factorial moment in the proposition, expand each falling factorial as an
ordered selection of distinct cycles.  Vertex-disjoint selections contribute the
product of their rooted-walk counts and the factors (16), giving
$\prod_\ell(\mu_\ell(1+\delta_\ell))^{b_\ell}$.  Every selection with an intersection
has a union with $e>v$ and is $O(n^{v-e})=O(1/n)$ after the same colour-compatible
matching calculation.  Since the total number and lengths are fixed, summing these
errors is $o(1)$.  This proves the asserted joint factorial moments.

It remains to identify the variance.  Since the spectral radius of $B$ is $3<4$ and
$\tr B=\tr B^2=0$,
\[
 \sum_{\ell\ge3}\mu_\ell\delta_\ell^2
 =2\sum_{\ell\ge1}\frac{\tr(B^\ell)}{\ell4^\ell}
 =-2\log\det(I-B/4).                                           \tag{17}
\]
For completeness, the Ihara--Bass identity here is
\[
 \det(I-uB)=(1-u^2)^5\det(I-uA+3u^2I).                         \tag{18}
\]
It follows by writing $B=XY-J$, where $J$ reverses directed edges and $X,Y$ are the
head and tail incidence matrices, then using $J^2=I$, $YX=A$, $YJX=4I$ and the
determinant lemma.  At $u=1/4$,
\[
 I-A/4+3I/16=(15/16)M,
\]
so $\det(I-B/4)=(15/16)^{10}\det M$.  Equations (17) and
\cref{prop:second} now give
\[
 \exp\left(\sum_{\ell\ge3}\mu_\ell\delta_\ell^2\right)
 =(16/15)^{20}\det(M)^{-2}=R.
\]
\end{proof}

\phantomsection\label{pf:main}
\thmmain*
\begin{proof}
We apply the small-subgraph conditioning theorem in the form stated in
\cite[Section~4]{Wormald99} to $Y_n=Z_n/\E Z_n$ and the cycle counts $X_{\ell,n}$.

\repair{The theorem's hypotheses are all explicit above: (i) finite collections of
$X_{\ell,n}$ converge to independent Poisson variables with means $\mu_\ell$ by
\cref{prop:cycles}; (ii) their mixed factorial moments under the $Z_n$-size-biased law
converge with factors $1+\delta_\ell$, again by \cref{prop:cycles};
(iii) $\delta_\ell>-1$ and
$\sum\mu_\ell\delta_\ell^2<\infty$; and (iv) the normalized second moment tends to
$\exp(\sum\mu_\ell\delta_\ell^2)$ by \cref{prop:second,prop:cycles}.  Therefore the
small-subgraph conditioning theorem applies directly, including the required exchange
of the limits in the number of conditioned cycle lengths and in $n$.}

It yields
\[
 Y_n\xrightarrow{d}W:=\prod_{\ell\ge3}
 (1+\delta_\ell)^{P_\ell}e^{-\mu_\ell\delta_\ell},              \tag{19}
\]
where the $P_\ell$ are independent Poisson variables of means $\mu_\ell$, and it gives
$\Pr(Z_n>0)\to1$.  To see positivity directly, the logarithm of the partial products is
\[
 \sum_{\ell=3}^L(P_\ell-\mu_\ell)\log(1+\delta_\ell)
 +\sum_{\ell=3}^L\mu_\ell(\log(1+\delta_\ell)-\delta_\ell).
\]
The first series converges almost surely because its variances are summable, and the
second converges absolutely; hence $W>0$ almost surely.  Thus $H_n$ is asymptotically
almost surely $3$-colourable.  Notice that the rounded $t_a=n/3+O(1)$ were used
throughout, so this holds for every integer sequence $n$, not only multiples of $3$.

We finally exclude chromatic number $2$.  In a proper $2$-colouring let $s_i$ be the
number of colour-$1$ vertices in fibre $i$.  Every base-edge matching forces
$s_i+s_j=n$.  Since $K_5$ contains triangles, $s_i=n/2$ for all $i$; hence no such
colouring exists for odd $n$.  For even $n$, the expected number of proper
$2$-colourings is exactly
\[
 \binom{n}{n/2}^{5}
 \left(\frac{((n/2)!)^2}{n!}\right)^{10}
 =\binom{n}{n/2}^{-5}=o(1).
\]
Markov's inequality shows that $H_n$ is asymptotically almost surely nonbipartite.
Together with $3$-colourability, this proves $\Pr(\chi(H_n)=3)\to1$.
\end{proof}

\section{Report of the LLM referee (verbatim)}\label{app:referee}

\begin{tcolorbox}[colback=gray!8,colframe=gray!50,boxrule=0.4pt,arc=2pt]
\small The following is the unedited report of the adversarial referee agent
(\texttt{claude-opus-5}, 2026-09-17) on the \emph{original} writeup. Its step numbers
refer to that writeup, not to this note. The scripts it mentions are in
\texttt{verification\_astra/scripts/chromatic\_number\_of\_random\_lifts\_of\_complete\_graphs/}.
Treat it as machine output, not as an authority.
\end{tcolorbox}

\input{.build/referee.tex}

\begin{thebibliography}{99}
\small
\setlength{\itemsep}{2pt}
\bibitem{OPG} A.~Amit, N.~Linial and J.~Matoušek, \emph{Chromatic number of random lifts of complete graphs}, OpenProblemGarden, posted 3 September 2012, \url{http://www.openproblemgarden.org/op/chromatic_number_of_random_lifts_of_complete_graphs}.
\bibitem{ALM02} A.~Amit, N.~Linial and J.~Matoušek, Random lifts of graphs: independence and chromatic number, \emph{Random Structures \& Algorithms} \textbf{20} (2002), 1--22.
\bibitem{FT12} B.~Farzad and D.~O. Theis, Random lifts of $K_5\setminus e$ are 3-colourable, \emph{SIAM Journal on Discrete Mathematics} \textbf{26} (2012), 169--176.
\bibitem{NPG21} J.~D. Nir and X.~Pérez-Giménez, The chromatic number of random lifts of complete graphs, arXiv:2109.13347 (2021).
\bibitem{Bregman73} L.~M. Bregman, Some properties of nonnegative matrices and their permanents, \emph{Soviet Mathematics Doklady} \textbf{14} (1973), 945--949.
\bibitem{Radhakrishnan97} J.~Radhakrishnan, An entropy proof of Bregman's theorem, \emph{Journal of Combinatorial Theory, Series A} \textbf{77} (1997), 161--164.
\bibitem{GJR10} C.~Greenhill, S.~Janson and A.~Ruciński, On the number of perfect matchings in random lifts, \emph{Combinatorics, Probability and Computing} \textbf{19} (2010), 791--817.
\bibitem{Wormald99} N.~C. Wormald, Models of random regular graphs, in \emph{Surveys in Combinatorics, 1999}, London Mathematical Society Lecture Note Series 267, Cambridge University Press, 1999, 239--298.
\end{thebibliography}

\end{document}
